\origpage{24}{}
Let $\Gamma, F_1^*, \ldots, F_r^*, H_1, \ldots, H_r, \mu$ be a presentation of $G$ in $\Sigma$, and let $L$ be a vital linkage in $G$. If $\Gamma \subseteq L$ and every terminal of $L$ is in one of $H_1, \ldots, H_r$, we say that $L$ is \textit{exhaustive} (for the presentation). If $L$ is exhaustive then it follows that $\Gamma$ is a forest and every vertex of $\Gamma$ not in $bd(\Sigma)$ has degree 2 in $\Gamma$. Let us say a $\Sigma$-\textit{jump} is a path in $\Gamma$ with distinct ends both in $bd(\Sigma)$ and with no internal vertex in $bd(\Sigma)$. If $L$ is exhaustive, then every edge of $\Gamma$ is in a unique $\Sigma$-jump, and any two $\Sigma$-jumps have no common vertices except possibly one end. The main result of this section is the following.

\textbf{6.3} \textit{For every surface $\Sigma$ and all integers $p, r \ge 0$ and $q \ge 1$, there exists $\lambda \ge 0$ such that, if $G$ is a graph with a presentation in $\Sigma$ of depth $\le q$ and defect $\le r$, and there is an exhaustive vital $p$-linkage in $G$, then there are at most $\lambda$ $\Sigma$-jumps.}

The most difficult step in the proof of (6.3) is where $\Sigma$ is a disc and $r = 1$. Then the cases when $\Sigma$ is a disc and $r = 2, 3$ and at least 4 are successively easier, and finally the case when $\Sigma$ is not a disc is also quite easy.

\textbf{6.4} \textit{For all integers $p \ge 0$ and $q \ge 1$ there exists $\lambda \ge 0$ such that, if $G$ has a presentation in a disc $\Sigma$ of depth $\le q$ and defect 1, and there is an exhaustive vital $p$-linkage in $G$, then there are at most $\lambda$ $\Sigma$-jumps.}

\textbf{Proof.} Choose $n$ so that (2.6) holds, with $k$ replaced by $2q$. Let $\lambda = n^n$. We claim that \del{(6.2)}\fix{(6.4)} is satisfied.\fixnote{fix:6.4-proof}{The reference meant is (6.4), since this is its proof. Result (6.2) has no $\lambda$ to choose, and the proof takes the hypotheses of (6.4) and ends with its conclusion.} For let $\Gamma, F^*, H, \mu$ be a presentation of $G$ in a disc $\Sigma$ with depth $\le q$ and let $L$ be a vital $p$-linkage which is exhaustive for the presentation.

A line $F \subseteq \Sigma$ is \textit{good} if it is proper, both its ends are in $V(\Gamma)$, and no internal point is in $U(\Gamma)$. If $F$ is a good line, there are two lines $J(F), K(F) \subseteq bd(\Sigma)$ with the same ends as $F$, where $K(F) \subseteq F^*$. Let $\Delta(F)$ be the closed disc in $\Sigma$ bounded by $F \cup K(F)$. Define
\begin{align*}
A(F) &= (\Gamma \cap \overline{\Sigma \setminus \Delta(F)}) \cup \mu(J(F)) \\
B(F) &= (\Gamma \cap \Delta(F)) \cup \mu(K(F)) \\
X(F) &= V(A(F) \cap B(F)).
\end{align*}
\hspace*{12pt}(1) \textit{(A(F), B(F)) is a separation of $G$, and $X(F) = V(\mu(u_1) \cup \mu(u_2))$ where $F$ has ends $u_1, u_2$.}

The proof is similar to that of theorem (2.1) of [10]\del{ and we omit it}. \fix{We give the proof. Let $W = V(B(F))$ and $X = X(F)$. Since $L$ is exhaustive, every edge of $\Gamma$ is in a unique $\Sigma$-jump, and every vertex of $\Gamma$ not in $bd(\Sigma)$ is an internal vertex of a $\Sigma$-jump. Since $F$ is good, $F \cap U(\Gamma) = \{u_1, u_2\} \subseteq bd(\Sigma)$, and so no $\Sigma$-jump meets $F$ except in its ends; hence a $\Sigma$-jump lies in $\Delta(F)$ if its ends are in $K(F)$, and in $\overline{\Sigma \setminus \Delta(F)}$ if they are in $J(F)$. In particular, $\Gamma \cap \Delta(F)$ is the union of the $\Sigma$-jumps with both ends in $K(F)$ and of the vertices of $\Gamma$ in $K(F)$, and a vertex of $\Gamma$ in both $\Delta(F)$ and $\overline{\Sigma \setminus \Delta(F)}$ is in $F$, and so is $u_1$ or $u_2$. Since $u_1, u_2 \in V(\Gamma)$, every space $s \subseteq F^*$ lies in $K(F)$ or in $J(F)$, and by (V3), $\mu(u) \subseteq \mu(s)$ when $u$ is an end of $s$. Since $u_1$ or $u_2$ lies in $F^*$ between any space in $K(F)$ and any space in $J(F)$, it follows from (V4) that}

\fix{\hspace*{12pt}(i) \textit{every vertex of $\mu(K(F)) \cap \mu(J(F))$ is in $X$, and no edge is in both $\mu(K(F))$ and $\mu(J(F))$.}}

\fix{Also $X \subseteq W \cap V(A(F))$, since $\mu(u_i) \subseteq \mu(K(F)) \cap \mu(J(F))$ ($i = 1, 2$). We claim that no edge $e$ of $G$ that is not in $B(F)$ has an end $x \in W \setminus X$. For suppose that $e$ is such an edge. If $e \in E(\Gamma)$, then $x$ is a vertex of $\Gamma \cap \Delta(F)$ other than $u_1, u_2$, so the $\Sigma$-jump that contains $e$ has $x$ as an internal vertex or as an end in $K(F)$, and so lies in $\Delta(F)$; but then $e \in E(B(F))$, a contradiction. Otherwise, by (V1) and (V2), $e$ is an edge of $\mu(s)$ for some space $s \subseteq F^*$, since the graphs $\mu(u)$ have no edges; and $s \subseteq J(F)$, since $e \notin E(B(F))$. If $x \in V(\Gamma)$, then $x$ is an end of $s$ by (V3), and $x \in K(F)$, so $x \in \{u_1, u_2\}$; and if $x \notin V(\Gamma)$, then $x \in V(\mu(K(F)))$, and $x \in X$ by (i). In both cases $x \in X$, a contradiction. This proves our claim.}

\fix{Consequently the edges of $G$ not in $B(F)$, and the vertices of $G$ not in $W \setminus X$, form a subgraph $A'$ of $G$, and $(A', B(F))$ is a separation of $G$ with $V(A' \cap B(F)) = X$. It remains to show that $A' = A(F)$. By (V1), (V2) and the above, every edge of $G$ is an edge of a $\Sigma$-jump or of some $\mu(s)$, and so belongs to $A(F)$ or to $B(F)$; and by (i) and since each $\Sigma$-jump lies on one side of $F$, no edge belongs to both. So $E(A') = E(A(F))$. Every vertex of $G$ not in $W$ is in $A(F)$, by (V1) and (V2). Finally, let $x \in V(A(F)) \cap W$. If $x \in V(\Gamma)$, then $x$ is a vertex of $\Gamma$ in both $\Delta(F)$ and $\overline{\Sigma \setminus \Delta(F)}$ (because $V(\mu(K(F)) \cap \Gamma) \subseteq K(F)$ and $V(\mu(J(F)) \cap \Gamma) \subseteq J(F)$, by (V3) and the definition of $\mu(u)$), and so $x \in \{u_1, u_2\} \subseteq X$; and if $x \notin V(\Gamma)$, then $x \in \mu(K(F)) \cap \mu(J(F))$, and $x \in X$ by (i). Hence $V(A') = V(A(F))$, and so $A' = A(F)$. This proves (1).}\fixnote{fix:6.4-claim1-proof-added}{The paper omits the proof of Claim (1). The argument above is the one used in the formalization.}
